\documentclass[11pt]{article}

\usepackage{amsmath, amssymb, amsthm}
\usepackage{geometry}
\usepackage{enumitem}
\usepackage{setspace}
\usepackage{hyperref}

\geometry{margin=1in}
\setstretch{1.2}

\title{\textbf{CSE400 -- Fundamentals of Probability in Computing}\\
Lecture 5: Bayes' Theorem, Random Variables, and Probability Mass Function}
\author{Instructor: Dhaval Patel, PhD}
\date{January 20, 2026}

\begin{document}

\maketitle

\section{Bayes' Theorem}

\subsection{Weighted Average of Conditional Probabilities}

Let $A$ and $B$ be events.

We may express event $A$ as:
\[
A = AB \cup AB^c
\]

This follows because, for an outcome to be in $A$, it must either:
\begin{itemize}[leftmargin=1.5em]
    \item be in both $A$ and $B$, or
    \item be in $A$ but not in $B$.
\end{itemize}

Since $AB$ and $AB^c$ are mutually exclusive, by \textbf{Axiom 3}:
\[
\Pr(A) = \Pr(AB) + \Pr(AB^c)
\]

Using conditional probability:
\[
\Pr(A) = \Pr(A \mid B)\Pr(B) + \Pr(A \mid B^c)[1 - \Pr(B)]
\]

\textbf{Interpretation:}  
The probability of event $A$ is a \textbf{weighted average of conditional probabilities}, where the weights are the probabilities of the events on which the conditioning is done.

\subsection{Learning by Example}

\subsubsection*{Example 3.1 (Part 1/2)}

An insurance company divides people into two classes:
\begin{itemize}[leftmargin=1.5em]
    \item accident prone
    \item not accident prone
\end{itemize}

Given:
\begin{itemize}[leftmargin=1.5em]
    \item $\Pr(\text{accident} \mid \text{accident prone}) = 0.4$
    \item $\Pr(\text{accident} \mid \text{not accident prone}) = 0.2$
    \item $\Pr(\text{accident prone}) = 0.3$
\end{itemize}

\textbf{Question:}  
What is the probability that a new policyholder will have an accident within one year?

\textbf{Solution:}

Let:
\begin{itemize}[leftmargin=1.5em]
    \item $A_1$: policyholder has an accident within one year
    \item $A$: policyholder is accident prone
\end{itemize}

\[
\Pr(A_1) = \Pr(A_1 \mid A)\Pr(A) + \Pr(A_1 \mid A^c)\Pr(A^c)
\]

\[
= (0.4)(0.3) + (0.2)(0.7) = 0.26
\]

\subsubsection*{Example 3.1 (Part 2/2)}

\textbf{Question:}  
Suppose a new policyholder has an accident within a year. What is the probability that they are accident prone?

\textbf{Solution:}

\[
\Pr(A \mid A_1) = \frac{\Pr(AA_1)}{\Pr(A_1)}
\]

\[
= \frac{\Pr(A)\Pr(A_1 \mid A)}{\Pr(A_1)}
= \frac{(0.3)(0.4)}{0.26} = \frac{6}{13}
\]

\subsection{Law of Total Probability}

Suppose $B_1, B_2, \ldots, B_n$ are mutually exclusive events such that:
\[
\bigcup_{i=1}^{n} B_i = B
\]

Exactly one of the events $B_1, B_2, \ldots, B_n$ must occur.

Writing:
\[
A = \bigcup_{i=1}^{n} AB_i
\]

Since the events $AB_i$ are mutually exclusive:
\[
\Pr(A) = \sum_{i=1}^{n} \Pr(AB_i)
= \sum_{i=1}^{n} \Pr(A \mid B_i)\Pr(B_i)
\]

This is called the \textbf{Law of Total Probability (Formula 3.4)}.

\subsection{Bayes Formula}

Using:
\[
\Pr(AB_i) = \Pr(B_i \mid A)\Pr(A)
\]

We obtain:
\[
\Pr(B_i \mid A) =
\frac{\Pr(A \mid B_i)\Pr(B_i)}
{\sum_{j=1}^{n} \Pr(A \mid B_j)\Pr(B_j)}
\]

This is known as the \textbf{Bayes Formula (Proposition 3.1)}.

\begin{itemize}[leftmargin=1.5em]
    \item $\Pr(B_i)$: \textbf{apriori probability}
    \item $\Pr(B_i \mid A)$: \textbf{posteriori probability}
\end{itemize}

\subsection{Learning by Example}

\subsubsection*{Example 3.2 (Card Problem)}

Three cards:
\begin{itemize}[leftmargin=1.5em]
    \item one red--red (RR)
    \item one black--black (BB)
    \item one red--black (RB)
\end{itemize}

A card is chosen at random and placed on the ground.  
The upper side is red.

\textbf{Question:}  
What is the probability that the other side is black?

\textbf{Solution:}

Let:
\begin{itemize}[leftmargin=1.5em]
    \item RR, BB, RB denote the chosen card type
    \item $R$: upturned side is red
\end{itemize}

\[
\Pr(RB \mid R) =
\frac{\Pr(R \mid RB)\Pr(RB)}
{\Pr(R \mid RR)\Pr(RR) + \Pr(R \mid RB)\Pr(RB) + \Pr(R \mid BB)\Pr(BB)}
\]

\[
= \frac{(1/2)(1/3)}
{(1)(1/3) + (1/2)(1/3) + (0)(1/3)}
= \frac{1}{3}
\]

\section{Random Variables}

\subsection{Motivation and Concept}

In many experiments, interest lies in a \textbf{function of the outcome}, not the outcome itself.

Examples discussed in lecture:
\begin{itemize}[leftmargin=1.5em]
    \item sum of two dice
    \item number of heads in coin tossing
\end{itemize}

A \textbf{random variable} is a real-valued function defined on the sample space.

\textbf{Definition:}  
A random variable $X$ on a sample space $\Omega$ is a function:
\[
X : \Omega \rightarrow \mathbb{R}
\]

Values are determined by outcomes of the experiment.  
Probabilities are assigned to possible values.

\subsection{Distribution of a Random Variable}

Two important components:
\begin{enumerate}[leftmargin=1.5em]
    \item the set of values the random variable can take
    \item the probabilities with which it takes those values
\end{enumerate}

For any value $a$:
\[
\{ \omega \in \Omega : X(\omega) = a \}
\]
is an event, abbreviated as $(X = a)$.

The collection of probabilities:
\[
\Pr(X = a)
\]
for all possible values $a$ is called the \textbf{distribution} of $X$.

\subsection{Visualization}

The distribution of a random variable can be visualized as a \textbf{bar diagram}:
\begin{itemize}[leftmargin=1.5em]
    \item x-axis: values of the random variable
    \item height of bar at $a$: $\Pr[X = a]$
\end{itemize}

\subsection{Example}

\textbf{Experiment:} Tossing 3 fair coins  

Let $Y$ = number of heads.

Possible values: $0, 1, 2, 3$

\[
\Pr(Y=0)=\frac{1}{8}, \quad
\Pr(Y=1)=\frac{3}{8}, \quad
\Pr(Y=2)=\frac{3}{8}, \quad
\Pr(Y=3)=\frac{1}{8}
\]

Since $Y$ must take one of these values:
\[
1 = \sum_{i=0}^{3} \Pr(Y=i)
\]

\section{Probability Mass Function (PMF)}

\subsection{Concept}

A random variable that can take \textbf{at most a countable number of possible values} is said to be \textbf{discrete}.

Let $X$ be a discrete random variable with range:
\[
R_X = \{x_1, x_2, x_3, \ldots\}
\]

The function:
\[
p(x) = \Pr(X = x)
\]
is called the \textbf{Probability Mass Function (PMF)} of $X$.

Since $X$ must take one of its possible values:
\[
\sum_x p(x) = 1
\]

\subsection{Example}

The PMF of a random variable $X$ is given by:
\[
p(i) = c \frac{\lambda^i}{i!}, \quad i = 0,1,2,\ldots
\]
where $\lambda > 0$.

Since:
\[
\sum_{i=0}^{\infty} p(i) = 1
\]

\[
c \sum_{i=0}^{\infty} \frac{\lambda^i}{i!} = 1
\]

Using:
\[
e^{\lambda} = \sum_{i=0}^{\infty} \frac{\lambda^i}{i!}
\]

\[
c e^{\lambda} = 1 \Rightarrow c = e^{-\lambda}
\]

Thus:
\[
\Pr(X=0) = e^{-\lambda}
\]

\[
\Pr(X>2) = 1 - \Pr(X \le 2)
= 1 - \Pr(X=0) - \Pr(X=1) - \Pr(X=2)
\]

\section{Class Participation}

Students are instructed to switch to \textbf{Campuswire} for participation and discussion.

\end{document}

